\section{Discussion, methodological caveats, and policy implications}
\label{sec:discussion}

\subsection{Why apparent Indonesian bitcoin segmentation reflects conversion frictions}
Studies comparing rupiah-quoted bitcoin directly against global tether prices often report slow adjustment or persistent premia and infer that domestic capital controls or retail biases isolate local crypto order books \citep{maulana2026segmentasi,makarov2020trading}. Separating currency and stablecoin pricing with synchronized five-minute venue-specific observations tells a narrower story: the pure crypto leg on Indodax behaves like an integrated satellite book bounded by statutory trading costs. Mean BTC pure-crypto deviation is $-0.83\text{ bps}$ ($95\%$ CI $[-2.26, +0.60]\text{ bps}$), $98.44\%$ of business-day observations sit within the dynamic domestic two-leg taker fee and tax band, and the local \texttt{USDT/IDR}--JISDOR stablecoin--FX wedge accounts for $83.10\%$ of total basis variance (or $82.39\%$ when volume-averaged over JISDOR's monitoring window).

\subsection{Economic constraints on arbitraging the stablecoin--FX basis}
Why does $B_t^{stablecoin\text{-}FX}$ average $+30.77\text{ bps}$ (SD $73.66\text{ bps}$, 2021--2026) instead of collapsing toward zero? Unlike Route A and Route B crypto legs, arbitraging the \texttt{USDT/IDR}--JISDOR wedge is not a symmetric on-exchange loop. Buying USDT abroad at wholesale USD parity and selling \texttt{USDT/IDR} locally requires cross-border wire settlement, Indonesian banking intermediation, foreign-exchange spread conversion away from the mid-market JISDOR fixing, domestic rupiah settlement delays, regulatory reporting compliance under OJK and Bank Indonesia FX monitoring, and seller final PPh 22 taxes ($21\text{ bps}$ under PMK 50/2025 plus platform and CFX charges). Even when the stablecoin--FX basis drops negative during onshore USDT liquidity surges (as seen prior to October 2026 in event windows), reversing the trade into physical USD interbank balances faces settlement friction and capital-transfer costs. Consequently, $B_t^{stablecoin\text{-}FX}$ reflects fiat-rail conversion friction and onshore stablecoin demand fluctuations rather than costless unexploited profit \citep{lyons2023stablecoins,lee2022kimchi}.

\subsection{Residual verification gaps and methodological limits}
Four empirical limitations should be weighed when interpreting these estimates:
\begin{enumerate}
  \item \textbf{Partial historical verification of mid-sample Indodax fee cuts:} While pre-2022 ($0/30\text{ bps}$) and late-2024 through November 2026 schedules (including the exact five-decimal October 2026 Pro rates) are backed by primary exchange disclosures and tax statutes, day-by-day historical Pro maker/taker tiers on Indodax USDT pairs between May 2022 and October 2023 remain partially documented publicly (\texttt{PARTIALLY\_VERIFIED} in Table~\ref{tab:fees}). Containment rates inside dynamic cost bands across 2022--2023 therefore depend slightly on whether $0/30$ or $10/20$ platform commissions applied to a specific user tier.
  \item \textbf{Single multi-year domestic venue for IDR time series:} Although we tested public APIs on Tokocrypto, Pintu, and Reku, none provided downloadable multi-year historical OHLCV series comparable to Indodax. Our cross-venue comparison against Binance's direct IDR pairs (available only from November 29, 2025 onward; $N = 100,362$ five-minute bars) shows close alignment ($-0.19\text{ bps}$ mean gap on \texttt{BTC/IDR}, $-2.11\text{ bps}$ on \texttt{USDT/IDR}), yet long-horizon inferences prior to late 2025 remain Indodax-specific.
  \item \textbf{Unidentifiable daily Cholesky information shares and residual diagnostics:} Daily bivariate VECMs exhibit cross-market residual correlation $\rho_u = 0.991$ because global crypto shocks hit local and offshore books within the same 24-hour window, making Hasbrouck Information Share bounds extremely wide ($[7.1\%, 98.2\%]$). Similarly, high-frequency five-minute VECMs leave residual autocorrelation and ARCH effects ($\text{LB } p < 0.001$, $\text{ARCH } p < 0.001$), motivating Newey--West HAC standard errors throughout rather than homoskedastic likelihood standard errors alone.
  \item \textbf{Absence of historical order-book snapshots and withdrawal latency data:} Realized order-book depth slippage ($4.35$--$9.74\text{ bps}$ for $\$1\text{k}$--$\$50\text{k}$ across Routes A and B, versus $63.14$--$68.53\text{ bps}$ for Route C) is measured from live November 13, 2026 limit-order books rather than a five-year L2 replay. Historical slippage during acute volatility episodes such as March 2023 USDC de-pegging or September 2024 exchange maintenance may have been higher than November 2026 steady-state depth implies.
\end{enumerate}

\subsection{Policy implications for Bank Indonesia, OJK, and CFX}
Three concrete regulatory insights emerge. First, policymakers evaluating onshore crypto market efficiency should monitor the domestic \texttt{USDT/IDR}--JISDOR basis rather than unadjusted \texttt{BTC/IDR} premia, because domestic stablecoin conversion is the primary locus of macro-FX spillover and dollar-demand stress. Second, CFX and KKI's two-stage $50\%$ fee reductions in March and October 2026 measurably compressed statutory domestic friction by $3.33\text{ bps}$ round-trip on IDR pairs ($65.44\text{ bps} \rightarrow 62.11\text{ bps}$ taker), bringing domestic two-leg execution closer to offshore levels, whereas PMK 50/2025 ($0\%$ buyer VAT, $21\text{ bps}$ seller PPh 22) shifted tax incidence without altering symmetric round-trip statutory costs. Third, taxing onshore sell transactions at $21\text{ bps}$ alongside a $100\text{ bps}$ statutory rate on offshore venues explains why domestic liquidity stays anchored locally even as cross-venue parity persists within fee bands.
